% -*- coding: utf-8 -*-
\begin{quote}
\small \textbf{Instrucciones.} Estos ejercicios usan sus propios predicados, así que se hacen a lápiz. Para practicar más la formalización, la evaluación de fórmulas y los contramodelos, usa la aplicación \appmundo{}, con su diccionario de formas, tamaños, colores y posiciones.
\end{quote}

\section*{Bloque 1: formalización de enunciados del lenguaje natural en LPO}
\textit{En este bloque practicamos la formalización de enunciados al
  lenguaje de la lógica de predicados.}

\begin{pregunta}[Enunciados categóricos y propiedades individuales]\label{e5-1}
Formaliza los siguientes enunciados utilizando los símbolos de predicado indicados entre paréntesis:
\begin{enumerate}
    \item Las manzanas y las naranjas son gustosas y nutritivas.  
    ($M(x)$: «$x$ es una manzana»; $N(x)$: «$x$ es una naranja»; $G(x)$: «$x$ es gustosa»; $U(x)$: «$x$ es nutritiva»).
    \item Hay alimentos que solo se pueden comer si han sido cocinados.  
    ($A(x)$: «$x$ es un alimento»; $C(x)$: «$x$ se puede comer»; $K(x)$: «$x$ ha sido cocinado»).
    \item Sin frenos, no hay ningún coche seguro.  
    ($F(x)$: «$x$ tiene frenos»; $C(x)$: «$x$ es un coche»; $S(x)$: «$x$ es seguro»).
    \item No todo el mundo es rico, culto y educado, ni todos los ricos son educados y cultos.  
    ($R(x)$: «$x$ es rico»; $C(x)$: «$x$ es culto»; $E(x)$: «$x$ es educado». Dominio: todas las personas).
    \item No todas las cosas compradas a bajo precio son delicadas y quebradizas.  
    ($C(x)$: «$x$ es una cosa»; $B(x)$: «$x$ ha sido comprada a bajo precio»; $L(x)$: «$x$ es delicada»; $Q(x)$: «$x$ es quebradiza»).
\end{enumerate}
\end{pregunta}

\begin{pregunta}[Relaciones, cuantificación e identidad]\label{e5-2}
Dado el predicado binario $A(x, y)$ («$x$ ama a $y$») y las constantes individuales $a$ (Antonio) y $b$ (Belén), formaliza los siguientes enunciados:
\begin{enumerate}
    \item Antonio ama a Belén pero Belén no ama a Antonio.
    \item Todo el mundo ama a Belén.
    \item Belén ama a alguien que también la ama a ella.
    \item Belén ama a alguien que no es Antonio (lo cual no implica que no ame a Antonio).
    \item Antonio siempre es correspondido en su amor.
    \item Belén ama a una sola persona que no es ella misma.
    \item Antonio ama exactamente a dos personas (véase la sección~\ref{sec:cuantificadores-numericos}).
    \item Todo el mundo se ama a sí mismo.
    \item Todo el que ama es correspondido.
    \item Todo el que ama es amado.
    \item Hay quien no es correspondido en su amor.
    \item Todo el mundo ama a alguien distinto de sí mismo.
\end{enumerate}
\end{pregunta}

\begin{pregunta}\label{e5-3}
Utilizando los predicados:
\begin{itemize}
    \item $A(x, y)$: «$x$ aprueba la asignatura $y$»
    \item $E(x)$: «$x$ estudia filosofía»
    \item $F(x, y)$: «la asignatura favorita de $x$ es $y$»
    \item $S(x)$: «$x$ es una asignatura»
\end{itemize}
y las constantes individuales $a$ (Antonio), $b$ (Belén) y $l$ (Lógica), formaliza los siguientes enunciados:
\begin{enumerate}
    \item La lógica no es la asignatura favorita de ningún estudiante de filosofía.
    \item Hay al menos dos estudiantes de filosofía distintos que tienen la misma asignatura favorita.
    \item La asignatura favorita de Antonio no es la misma que la de Belén.
    \item Antonio y Belén aprueban exactamente las mismas asignaturas.
    \item Todos los estudiantes de filosofía cuya asignatura favorita es la lógica, la aprueban.
    \item Belén no aprueba la asignatura favorita de Antonio.
    \item Belén aprueba todas las asignaturas que Antonio suspende (entiende que suspender es no aprobar).
    \item La lógica no es la asignatura favorita de ningún estudiante de filosofía que la suspenda.
\end{enumerate}
\end{pregunta}

\begin{pregunta}[Anidamiento de cuantificadores]\label{e5-anid}
Formaliza los siguientes enunciados, indicando el diccionario:
\begin{enumerate}
    \item Todo filósofo ha leído a algún clásico.
    \item Hay un clásico que todos los filósofos han leído.
    \item Ningún estudiante aprueba todas las asignaturas.
    \item Toda asignatura tiene algún estudiante que no la aprueba.
    \item Hay un profesor al que no entiende ninguno de sus alumnos.
\end{enumerate}
¿Cuál de los enunciados 1 y 2 implica al otro? Da una estructura que muestre que el otro no lo implica.
\end{pregunta}

\begin{pregunta}[Cantidad y descripciones\ampliacion]\label{e5-cant}
Formaliza sin usar functores ni los operadores $\iota$ y $\varepsilon$:
\begin{enumerate}
    \item Hay al menos tres filósofos presocráticos.
    \item Belén suspende como mucho una asignatura.
    \item Hay exactamente dos asignaturas que Belén no aprueba.
    \item El autor de la \emph{Ética a Nicómaco} fue discípulo de Platón.
    \item La montaña de oro no existe (según el análisis de Russell).
\end{enumerate}
\end{pregunta}

\section*{Bloque 2: evaluación de fórmulas y verdad en estructuras}

\begin{pregunta}[Diseño de modelos]\label{e5-4}
Considera las oraciones formalizadas en el ejercicio~\ref{e5-2} sobre el predicado $A(x, y)$ («amar»). Diseña una estructura concreta $\mathcal{M} = \langle D, I \rangle$ con un dominio finito de cuatro personas $D = \{1, 2, 3, 4\}$ con $I(a) = 1$ e $I(b) = 2$ donde se verifiquen simultáneamente:
\begin{enumerate}
    \item Belén ama a alguien que también la ama a ella ($\exists x (A(b, x) \land A(x, b))$).
    \item Antonio siempre es correspondido en su amor ($\forall x (A(a, x) \to A(x, a))$).
    \item Hay quien no es correspondido en su amor ($\exists x \exists y (A(x, y) \land \neg A(y, x))$).
\end{enumerate}
Define explícitamente el conjunto de pares $I(A) \subseteq D \times D$ y justifica por qué satisface las tres fórmulas.
\end{pregunta}

\begin{pregunta}[Justificación semántica del silogismo clásico]\label{e5-5}
Demuestra formalmente, utilizando únicamente las definiciones semánticas de verdad y modelo, que el siguiente argumento es válido:
\begin{quote}
\textit{Los caracoles son gasterópodos.}\\
\textit{Los gasterópodos son moluscos.}\\
\textit{Por tanto, los caracoles son moluscos.}
\end{quote}
\textit{Pista.} Formaliza las premisas como $\forall x (C(x) \to
G(x))$ y $\forall x (G(x) \to M(x))$, y la conclusión como $\forall x
(C(x) \to M(x))$. Supón una estructura arbitraria $\mathcal{M} =
\langle D, I \rangle$ que hace verdaderas las premisas y razona para
cualquier elemento $d \in D$ observando que $I(C) \subseteq I(G)
\subseteq I(M)$. 
\end{pregunta}

\begin{pregunta}[Evaluación en una estructura aritmética infinita]\label{e5-6}
Sea la estructura $\mathcal{M} = \langle D, I \rangle$, donde $D = \mathbb{N} = \{0, 1, 2, 3, \dots\}$, $I(P)$ son los números pares, $I(Q)$ los números primos, e $I(R) = \le$. Determina si son verdaderas o falsas en $\mathcal{M}$:
\begin{enumerate}
    \item $\forall x \exists y R(x, y)$
    \item $\exists y \forall x R(y, x)$
    \item $\forall x \forall y (R(x, y) \lor R(y, x))$
    \item $\forall x (P(x) \to \exists y (Q(y) \land R(x, y)))$
    \item $\forall x \forall y ((R(x, y) \land x \neq y) \to \exists z (R(x, z) \land R(z, y) \land z \neq x \land z \neq y))$
\end{enumerate}
\end{pregunta}

\section*{Bloque 3: construcción de contramodelos y comparación semántica}

\begin{pregunta}[Diferenciación semántica de prefijos de cuantificadores]\label{e5-7}
Considera las siguientes cinco sentencias sobre una relación binaria $R$:
\begin{enumerate}
    \item $\exists x \exists y R(x, y)$
    \item $\forall x \forall y R(x, y)$
    \item $\exists x \forall y R(x, y)$
    \item $\forall x \exists y R(x, y)$
    \item $\exists y \forall x R(x, y)$
\end{enumerate}
Proporciona estructuras con dominio $D =
\{1, 2\}$ que muestren que ninguna de ellas es equivalente a otra (por ejemplo,
una estructura donde la fórmula 4 sea verdadera pero la 3 sea
falsa). ¿Hay alguna(s) que sea consecuencia lógica de otra(s)?\\
\textit{Pista.} Presta atención a la diferencia entre $\forall x \exists y R(x, y)$ (donde para cada $x$ puede elegirse un $y$ potencialmente distinto) y $\exists y \forall x R(x, y)$ (donde un mismo $y$ debe servir para todos los $x$).
\end{pregunta}

\begin{pregunta}[Refutación de falacias formales mediante contramodelos]\label{e5-8}
Construye explícitamente estructuras $\mathcal{M} = \langle D, I \rangle$ que hagan verdaderas todas las premisas y falsa la conclusión para refutar:
\begin{enumerate}
    \item $\forall x (P(x) \lor Q(x)) \ \not\models \ \forall x P(x) \lor \forall x Q(x)$
    \item $\forall x (P(x) \to Q(x)), \ \exists x \neg P(x) \ \not\models \ \exists x \neg Q(x)$
    \item $\forall x (P(x) \to Q(x)), \ \exists x Q(x) \ \not\models \ \exists x P(x)$
    \item $\forall x (P(x) \to Q(x)), \ \neg P(a) \ \not\models \ \neg Q(a)$
\end{enumerate}
\end{pregunta}

\section*{Bloque 4: requisitos de cardinalidad del dominio}

\begin{pregunta}[Condiciones de tamaño sobre el universo de discurso]\label{e5-9}
Justifica qué requisitos impone cada una de las siguientes fórmulas a la cardinalidad del dominio $D$ para que una estructura $\mathcal{M} = \langle D, I \rangle$ pueda satisfacerlas:
\begin{enumerate}
    \item $\exists x \exists y (x \neq y)$
    \item $\exists x \exists y (x \neq y \land \forall z (z = x \lor z = y))$
    \item $\forall x \exists y R(x, y) \ \land \ \forall x \forall y (R(x, y) \to \neg R(y, x)) \ \land \ \forall x \forall y \forall z ((R(x, y) \land R(y, z)) \to R(x, z))$
\end{enumerate}
\end{pregunta}

\begin{pregunta}[Estructuras de cardinalidad mínima para conjuntos de fórmulas]\label{e5-10}
Para cada uno de los siguientes tres conjuntos de sentencias de primer orden, propón una estructura $\mathcal{M} = \langle D, I \rangle$ que lo satisfaga procurando que el dominio tenga la menor cardinalidad posible. Justifica en cada caso por qué no es posible hacerlo con un número menor de elementos:
\begin{enumerate}
    \item $\Gamma_1 = \{P(a), \ P(b), \ P(c), \ P(d)\}$.  
    \textit{Pista.} Observa si el lenguaje exige explícitamente que las constantes se refieran a individuos distintos o si pueden coincidir.
    \item $\Gamma_2 = \{\forall x (\neg P(x) \to R(x, x)), \ \forall x (P(x) \leftrightarrow \neg Q(x)), \ P(a), \ \exists x Q(x)\}$.
    \item $\Gamma_3 = \{\exists x \exists y \exists z (x \neq y \land y \neq z \land x \neq z), \ \forall x R(x, x), \ \forall x \forall y (R(x, y) \to R(y, x)), \ \forall x \forall y \forall z ((R(x, y) \land R(y, z)) \to R(x, z))\}$.
\end{enumerate}
\end{pregunta}

\section*{Bloque 5: importación existencial}

\begin{pregunta}[Darapti]\label{e5-11}
Considera el modo silogístico Darapti: «todo $M$ es $P$; todo $M$ es $S$; por tanto, algún $S$ es $P$».
\begin{enumerate}
    \item Construye una estructura que muestre que $\forall x (M(x) \to P(x)), \ \forall x (M(x) \to S(x)) \not\models \exists x (S(x) \land P(x))$.
    \item Demuestra en deducción natural que el argumento es válido si se añade la premisa $\exists x M(x)$.
\end{enumerate}
\end{pregunta}

\begin{pregunta}[El cuadro de oposición]\label{e5-12}
Sean A: $\forall x (S(x) \to P(x))$, E: $\forall x (S(x) \to \neg P(x))$, I: $\exists x (S(x) \land P(x))$ y O: $\exists x (S(x) \land \neg P(x))$.
\begin{enumerate}
    \item Muestra que A y O no pueden ser ambas verdaderas ni ambas falsas en ninguna estructura (son contradictorias). ¿Ocurre lo mismo con E e I?
    \item Construye una estructura en la que A y E sean ambas verdaderas, y otra en la que I y O sean ambas falsas.
    \item ¿Qué premisa habría que añadir para recuperar la subalternación (de A a I)?
\end{enumerate}
\end{pregunta}

%%% Local Variables:
%%% mode: latex
%%% TeX-master: "ApuntesLogica"
%%% End:
